\documentclass{situationist-preprint}

\title{The Irreducible Components of the Good}
\msid{ETH-01-b}
\msfamily{ETH-01}
\msdiscipline{Normative and Applied Ethics}
\mstype{Companion manuscript}
\mskeywords{the good; value pluralism; objective list; irreducibility; supervenience; method of isolation; well-being; knowledge; achievement; relationships; fairness; beauty; autonomy}

\begin{document}
\maketitle

\begin{abstract}
ETH-01 holds that the good is what idealized reflective endorsement converges on and that it has seven irreducible components: well-being, knowledge, achievement, relationships, fairness, beauty and autonomy. Volume~I of \emph{On the Good} (ETH-01-a) proves plurality from seven isolation verdicts. This companion examines what the irreducibility of a component is and then takes the components one by one. We distinguish three notions of irreducibility and show that the formal one, failure of supervenience on the other six, is equivalent to membership in every supervenience base of value, so that, given the completeness thesis of Volume~VI, the seven form the unique minimal base. We show that irreducibility implies neither separability nor incommensurability, and that the isolation test compares levels of components, so that its pairs may be built by compensation. For each component we state the isolation pair, the most serious attempt to reduce it to the others, and the reason the attempt fails.
\end{abstract}

\section{Introduction}\label{s:intro}

\subsection{The claim}

The conclusion of ETH-01 has two parts. The first is a criterion: the good is what idealized reflective endorsement converges on, so that one outcome is better than another if and only if every agent's idealized reflective endorsement would prefer it (Definition~3.2 of ETH-01-a). The second is a result about what satisfies the criterion: the good has seven components, each irreducible to the others. Volume~I establishes the second part in a compact form. It represents an outcome by a profile that assigns a level to each of the seven components, takes as its single substantive premise the conjunction of seven isolation verdicts (Assumption~5.1), and proves that the value of outcomes does not supervene on any set of components that omits one of the seven (Theorem~5.2). Hedonism and every other monism follow as casualties (Corollaries~5.3 and~5.4).

That proof is short, and its brevity is instructive: as ETH-01-a remarks, the whole weight lies on the isolation verdicts. The present manuscript is about those verdicts and about the notion of irreducibility they support. It has two aims. The first is conceptual. \enquote{Irreducible} is used in several senses in the literature on value pluralism, and some of the objections to objective lists trade on the difference between them. We fix the sense in which ETH-01 asserts irreducibility, relate it to the others, and draw out two consequences that are often denied: that irreducible components may interact, and that they may be fully comparable. The second aim is to review the components in turn. For each we give the pair of outcomes that shows it passes the isolation test of Definition~4.1 of ETH-01-a, the most serious attempt to reduce it to the other six, and the reason the attempt fails. The extended treatment of each component belongs to Volumes~II to V of ETH-01; the present review states the core of each case in a form that can be assessed on its own.

\subsection{The charge of arbitrariness}

Objective-list theories have long been charged with being lists and nothing more. \citet[ch.~3]{Sumner1996} argued that such theories enumerate goods without explaining what makes them good, and that their lists are therefore at the mercy of whichever intuitions their authors happen to have. \citet{Fletcher2013} answered that an objective-list theory need not supply a single explanatory property, since its explanatory work can be done item by item, and \citet{Rice2013} argued that the demand for a unifying explanation begs the question against pluralism. These replies concern the list of prudential goods, but the charge applies with equal force to a list of the goods of outcomes, and it has two parts. One is that the list may be too long: some of its items may be reducible to others, so that the list counts a single good twice. The other is that it may be too short: something not on it may matter. ETH-01-d answers the second part by identifying the completeness of the list with the supervenience of value on all seven components, the thesis of Volume~VI. The present manuscript answers the first part. A list counts a good twice only if one of its items is reducible to the others, and the isolation verdicts are precisely the claims that none is.

\subsection{Plan}

Section~\ref{s:notions} distinguishes three notions of irreducibility and proves the structural results. Section~\ref{s:method} sets out how isolation pairs are constructed. Section~\ref{s:components} takes the seven components in turn. Section~\ref{s:composite} considers composite and global reductions. Section~\ref{s:obj} answers objections, and Section~\ref{s:not} states what the discussion does not establish.

\section{Three notions of irreducibility}\label{s:notions}

We use the framework of ETH-01-a. An outcome is represented by a profile $x$ assigning to each component $i$ among the seven a level $x_i$, and $V$ is the value assigned to profiles by idealized reflective endorsement. We write $\mathcal{C}$ for the set of seven components. For a set $S \subseteq \mathcal{C}$, $V$ \emph{supervenes} on $S$ if any two profiles that agree on every component in $S$ have the same value; we call such a set a \emph{supervenience base} for $V$.

\begin{definition}\label{d:formal}
A component $i$ is \emph{formally irreducible} if $V$ does not supervene on $\mathcal{C}\setminus\{i\}$.
\end{definition}

\begin{definition}\label{d:essential}
A component $i$ is \emph{base-essential} if it belongs to every supervenience base for $V$.
\end{definition}

\begin{definition}\label{d:nonderiv}
A component $i$ is \emph{non-derivative} if the contribution of $i$ to the value of outcomes is not explained by its contribution to the levels of the other components.
\end{definition}

Formal irreducibility is the notion of Theorem~5.2 of ETH-01-a: by Assumption~5.1, every component is formally irreducible. Non-derivativeness is the notion that the debate about instrumental and final value has in view. Knowledge is derivative, in this sense, if it is good only because it makes people's lives go better or enables them to achieve things; fairness is derivative if it is good only because fair procedures tend to produce outcomes that are better in well-being. The two notions are related but not the same. A component whose value is wholly derivative cannot pass the isolation test, since if its contribution is exhausted by its effects on the other components, two outcomes alike in those components do not differ in value. A component that passes the isolation test may still have derivative value in addition: knowledge both is good and makes lives go better. Formal irreducibility is therefore the claim that each component has some non-derivative value, not that its value is wholly non-derivative.

\begin{proposition}\label{p:essential}
A component is formally irreducible if and only if it is base-essential.
\end{proposition}

\begin{proof}
Supervenience is preserved under enlarging the base: if $V$ supervenes on $S$ and $S \subseteq T$, any two profiles that agree on $T$ agree on $S$ and so have the same value. Suppose $i$ is formally irreducible and $S$ is a supervenience base that omits $i$. Then $S \subseteq \mathcal{C}\setminus\{i\}$, so $V$ supervenes on $\mathcal{C}\setminus\{i\}$, which contradicts formal irreducibility. Conversely, if $i$ is base-essential, then $\mathcal{C}\setminus\{i\}$, which omits $i$, is not a supervenience base, and $i$ is formally irreducible.
\end{proof}

Proposition~\ref{p:essential} reformulates Theorem~5.2 of ETH-01-a: every supervenience base for $V$ contains all seven components. It also clarifies what the theorem does not say. It does not say that any base exists among the subsets of $\mathcal{C}$; that is the completeness thesis.

\begin{corollary}\label{c:unique}
If the list is complete, so that $V$ supervenes on $\mathcal{C}$, then $\mathcal{C}$ is the unique minimal supervenience base for $V$ among sets of components.
\end{corollary}

\begin{proof}
By completeness $\mathcal{C}$ is a base. By Proposition~\ref{p:essential} and Theorem~5.2 of ETH-01-a, every base contains every member of $\mathcal{C}$, and so is $\mathcal{C}$ itself. Hence $\mathcal{C}$ is minimal, and no other set of components is a base at all.
\end{proof}

ETH-01-d shows that, given completeness, the seven are a minimal base (its Corollary~1.2). Corollary~\ref{c:unique} adds uniqueness. This matters for the charge of arbitrariness. If the components could be chosen in more than one way, as the axes of a space can be rotated, the list would reflect a choice of coordinates, and its items would have no better claim to be the components of the good than any other system of coordinates. Uniqueness shows that, among sets of the seven candidate dimensions, there is no such latitude. It does not exclude a redescription of the same outcomes by quite different dimensions, which we take up in Section~\ref{s:composite}.

Two further consequences concern what irreducibility does not imply. Say that $V$ is \emph{separable} in a component $i$ if the ranking of two profiles that agree on $i$ never depends on the level of $i$ at which they agree, and that two components are \emph{commensurable} if a gain in one can be outweighed by a loss in the other.

\begin{proposition}\label{p:interact}
Formal irreducibility of all seven components is consistent with the failure of separability in each of them, and with the full commensurability of every pair of them.
\end{proposition}

\begin{proof}
Let the levels be non-negative integers and let $V(x) = (x_1 + 1)(x_2 + 1)\cdots(x_7 + 1)$. Raising any single component by one raises $V$, so for each component there are two profiles that agree on the other six and differ in value, and each component is formally irreducible. $V$ is not separable in any component: the gain in value from raising component $j$ by one is $V(x)/(x_j+1)$, which depends on the level of every other component. And every pair is commensurable, since $V$ is real-valued: a profile with a higher level of $i$ and a lower level of $j$ than another is better, worse or equally good, according to the numbers.
\end{proof}

The function in the proof is chosen for simplicity, not as a proposal about the shape of $V$. Its point is that the two theses often bundled with pluralism are independent of it. Irreducibility is a claim about what value depends on, not about how it depends on it. That the value of fairness may depend on the level of well-being, that a beautiful world with no one in it may be less valuable than a beautiful world with people to see it, or that knowledge of trivial facts adds little, are claims about interaction, and they are compatible with the formal irreducibility of each component. This agrees with the recognition, in the literature on intrinsic value, that the value a thing has for its own sake may be conditional on its relations to other things \citep{Korsgaard1983,Kagan1998,RabinowiczRonnow2000}. Whether some pairs of components are in fact incommensurable, or on a par in the sense of \citet{Chang2002}, is the subject of Chapter~6 of Volume~I; Proposition~\ref{p:interact} shows only that irreducibility does not settle it.

\section{How isolation pairs are built}\label{s:method}

The isolation test asks for two outcomes that \enquote{differ only in their level of} a component and that idealized reflective endorsement ranks differently (Definition~4.1 of ETH-01-a). Two features of this formulation are easily overlooked, and both matter for the cases below.

The first is that the test compares levels of components, not concrete outcomes. Two outcomes that differ only in their level of knowledge may differ in a great many concrete respects: who knows what, which experiences people have, which events occur. What is required is that the differences cancel at the level of the six other components. This permits \emph{construction by compensation}. Suppose a case designed to isolate relationships, in which a friendship is genuine in one outcome and feigned in the other, inevitably also differs in knowledge, since the deceived party holds a false belief that the undeceived party knows to be true. The pair can be completed by giving the deceived party, in the second outcome, an item of knowledge of equal significance elsewhere, so that the two outcomes agree in level of knowledge while differing in relationships. The verdict about the completed pair is the one that tests relationships.

\begin{remark}\label{r:comp}
Construction by compensation presupposes that the levels of the components are measured on scales that allow equal differences to be identified, at least locally. It does not presuppose that the components are commensurable with one another, since compensation takes place within a single component.
\end{remark}

The second feature is that the test requires idealized reflective endorsement to rank the two outcomes differently, not that it ranks them by a large margin. A component may pass with a weak verdict. This matters for beauty, where the verdict is weaker than for the other components, as ETH-01-a acknowledges, and for autonomy, where the verdict depends on a careful statement of the case. A weak but convergent verdict suffices for irreducibility; it bears only on the weight of the component, which is a separate question.

A third point concerns the judge. The idealized endorser of Definition~3.1 of ETH-01-a is fully informed about the outcomes compared and their alternatives, imagines them vividly, reasons without error, and considers them from a standpoint that gives equal weight to everyone affected. Three common sources of error in isolation judgments are thereby excluded: the importation of unstated differences, which full information about the outcomes forbids; the substitution of a choice for a comparison, which the impartial standpoint forbids, since the judge is not choosing for herself; and the distortion that comes from imagining one outcome more vividly than the other. These are the sources of error that the debunking literature on the experience machine identifies \citep{DeBrigard2010,Weijers2014}, and ETH-05-a shows how a comparative statement of a case avoids them. The pairs below are stated in the comparative form.

\section{The seven components}\label{s:components}

For each component we give the isolation pair, the principal reduction, and the reason the reduction fails. The pairs are stated in outline; the full statements, with the compensations made explicit, are those of Chapter~4 of Volume~I and of the volume devoted to the component.

\subsection{Well-being}

\paragraph{The pair.} Two outcomes are alike in knowledge, achievement, relationships, fairness, beauty and autonomy. In the first, a person who will never know or remember it suffers a severe headache for an hour; in the second, she does not. The second is better.

\paragraph{The reduction.} No one proposes to reduce well-being to the other six components in the way that knowledge or fairness has been reduced to well-being. The live proposal runs the other way, and it concerns the relation between well-being and the other components rather than the irreducibility of well-being itself. Objective-list theories of well-being \citep{Parfit1984,Griffin1986,Fletcher2013} hold that knowledge, achievement and relationships are constituents of a person's well-being. If they are, it may seem that a profile which records well-being separately counts them twice.

\paragraph{Why it fails.} The pair shows that well-being has a dimension, the experiential, that varies independently of the other six. ETH-05-a adds a second dimension: a life in contact with reality goes better for its subject than an experiential twin that is not (its Theorem~3.2 and Corollary~3.3). The profile's well-being component records how a person's life goes for her in these respects. Whether knowledge, achievement and relationships are, in addition, constituents of well-being is not a question the profile needs to settle. If they are, the profile records the part of their value that is a person's good under the heading of the component itself, and the well-being level is measured net of it, as Corollary~5.4 of ETH-01-a allows. Either way, nothing is counted twice, since the isolation test for each component holds the others fixed. The constituent view does bear on one question: whether fairness, beauty and autonomy could also be relocated inside well-being. Chapter~7 of Volume~I and Theorem~10.4 of Volume~IV (ETH-01-e, E5) argue that, for fairness at least, they cannot, and we return to beauty and autonomy below.

\subsection{Knowledge}

\paragraph{The pair.} Two outcomes are alike in the six other components. In the first, the members of a community know the history of their ancestors: their beliefs about it are true and are held on the basis of good evidence transmitted reliably. In the second, they hold the same true beliefs, but the beliefs rest on a forged chronicle that happens to be accurate, and the members are compensated, as Section~\ref{s:method} describes, by an item of knowledge of equal significance in some other domain. No experience, achievement or relationship differs. The first outcome is better. A second pair isolates the difference between ignorance and knowledge rather than between true belief and knowledge: two outcomes alike in everything else, in one of which astronomers come to know the composition of a distant galaxy, a fact that will make no practical difference and will give no one more pleasure than the alternative belief, and in the other of which they do not. The first is better.

\paragraph{The reductions.} Two reductions have been proposed. The first holds that knowledge is valuable only because true belief is useful and knowledge is a reliable source of true belief; on this view the value of knowledge reduces to the value of the well-being it serves. The second, associated with virtue epistemology, holds that knowledge is a cognitive achievement, belief that is true because of the exercise of the believer's competence \citep{Sosa2007,Greco2010}. If that is right, the value of knowledge may seem to reduce to the value of achievement.

\paragraph{Why they fail.} The first pair answers the instrumental reduction directly, since the two outcomes agree in what people believe and in what they can do. It also answers the version of the reduction on which knowledge is worth more than true belief only because it is more useful: by EPI-14, knowledge is more valuable than mere true belief because it is more stable, and stability, though it is often useful, is a feature of a person's cognitive relation to the world, which the second outcome lacks whatever its effects. The second pair answers the reduction to achievement. The astronomers' knowledge is an achievement, but the knowledge of the community in the first pair is mostly not: it is acquired by testimony, and testimonial knowledge is not, in the general case, creditable to the hearer \citep{Lackey2007}. EPI-08 holds that testimony transmits knowledge without generating it; the hearer's knowledge is real, and the hearer has achieved nothing. Knowledge without achievement still counts in the comparison, and so knowledge is not the achievement component under another name. \citet{Kvanvig2003} has argued that the distinctive value lies in understanding rather than in knowledge; if so, the component should be read as including understanding, and the pairs are unaffected.

\subsection{Achievement}

\paragraph{The pair.} Two engineers spend their careers designing a bridge, with the same skill, effort and experiences. Neither lives to see the result. In the first outcome, the bridge stands and serves; in the second, a flaw in the steel supplied, which no inspection could have detected, causes it to be abandoned before it opens, and the users it would have served are served equally well by a ferry. Each engineer believes, with good evidence, that her bridge will stand; the second engineer's false belief is compensated by an item of knowledge of equal significance. The first outcome is better.

\paragraph{The reduction.} \citet{Bradford2015} analyzes achievement as the competent causation of a difficult outcome. On a natural reduction, the value of achievement lies in the knowledge exercised and the effort expended, both of which are present in the two outcomes, and in the well-being of those affected by the result, which the ferry equalizes.

\paragraph{Why it fails.} The pair holds the knowledge, the effort and the effects on others fixed, and the verdict survives. What differs is whether the engineer's competent effort issued in what it aimed at. This is what achievement adds, and it is not a matter of how anyone's life feels or of what anyone knows. The case also shows why achievement is not reducible to well-being on any mental-state theory, since the engineers' experiences are identical, and it agrees with ETH-05-a, on which a life whose aims are achieved goes better for its subject than an experiential twin whose aims are not. Achievement is a component of the good, not a condition of a meaningful life: EXI-17 argues that meaning requires active engagement with something worthwhile but not achievement, and the two claims are consistent, since a life may be meaningful without containing an achievement whose presence would make the outcome better.

\subsection{Relationships}

\paragraph{The pair.} Two outcomes are alike in the six other components. In the first, two people are friends: each cares about the other for the other's sake, and each knows it. In the second, each acts toward the other exactly as in the first, but out of a calculated interest in the other's patronage; neither ever discovers this, and the false belief of each that the other cares is compensated as before. Their experiences of the friendship are the same. The first outcome is better.

\paragraph{The reductions.} The value of friendship has been reduced to the well-being of the friends, either directly, since friendship is pleasant, or indirectly, since friends help one another. It has also been reduced to knowledge, on the view that what matters in a relationship is mutual understanding.

\paragraph{Why they fail.} The pair holds well-being in its experiential dimension fixed. It holds knowledge fixed by compensation. What differs is the attitude of each toward the other, the fact that makes the relationship one of friendship. \citet{Kolodny2003} argues that love is a form of valuing a relationship, and \citet{Scheffler1997} that relationships ground reasons that are not reducible to the value of the goods relationships produce; both views explain the verdict. ETH-19 holds that partiality toward those close to us is justified without impartial reasons, and the present verdict is the impartial counterpart: from a standpoint that gives equal weight to everyone, an outcome containing genuine relationships is better than one containing their behavioral replicas. The relationship component is not a component of anyone's well-being alone, since a friendship is a relation between two lives; whether the relationship also makes each life go better for the one who lives it is a question for Volume~III, and the irreducibility of the component does not depend on its answer.

\subsection{Fairness}

\paragraph{The pair.} Two outcomes contain the same total and the same distribution of well-being, knowledge, achievement, relationships, beauty and autonomy. In the first, the distribution results from a lottery to which everyone consented; in the second, from a procedure that favored some arbitrarily. The first is better. This is the pair of Chapter~4 of Volume~I and of Chapter~2 of Volume~IV, in the corrected form given in ETH-01-e (E1), where the agreement in the distribution of the other components is stated explicitly.

\paragraph{The reductions.} Fairness has been reduced to equality, to the priority of the worse off, and to desert. \citet{Diamond1967} pressed the point that a social ranking indifferent between giving a good to one person for certain and giving each of two persons an equal chance of it ignores the fairness of the procedure, and \citet{Broome1990} gave the matching account of fairness as the proportional satisfaction of claims.

\paragraph{Why they fail.} The corrected pair holds the distribution fixed, so a preference for the first outcome is not a preference for equality of outcome; and it holds the levels of the worse off fixed, so it is not a preference for priority. That fairness differs from equality agrees with POL-08, on which equality is not valuable in itself and the levelling-down objection \citep{Parfit1997} is decisive against egalitarianism: the levelling-down objection tells against a value that could be served by making the better off worse off, and fairness, on the corrected pair, is not such a value. Desert is not at issue in the pair, since no one in it deserves more than another. Volume~IV examines the three reductions in its Chapters~3, 6 and~9, and its Theorem~10.4 shows that fairness is not a component of anyone's well-being, since two outcomes alike in the well-being of every individual can differ in fairness and in value. Fairness is the clearest case of a good of outcomes that is not a good for anyone, and its presence on the list is the main respect in which the list differs from lists of prudential goods.

\subsection{Beauty}

\paragraph{The pair.} Two outcomes are alike in every respect, including the experiences of every sentient being, except that in one a region no one will ever see contains a beautiful formation and in the other a formless one. The first is better. This is the contrastive form of \citet[\S50]{Moore1903}, who asked us to compare a world exceedingly beautiful with one that is the ugliest we can conceive, neither ever to be seen by anyone.

\paragraph{The reduction.} \citet[III.xiv]{Sidgwick1907} replied that no one would think it rational to produce beauty in external nature apart from any possible contemplation of it, and that the value of beauty therefore lies in the experience of it. On this view, beauty reduces to well-being, specifically to the aesthetic experiences of those who perceive it.

\paragraph{Why it fails.} The pair holds every experience fixed, so if beauty mattered only through aesthetic experience, the two outcomes would be equally good. The verdict that the first is better is weaker than the verdicts for the other components, and ETH-01-a acknowledges that it is. But it is a verdict, and a weak verdict suffices for the isolation test (Section~\ref{s:method}). Sidgwick's argument concerns what it is rational to aim at, and the isolation test concerns which of two outcomes is better. The two come apart if, as is plausible, the value of unseen beauty is small enough that it would rarely be rational to sacrifice anything for it. Moore himself, in \S113, held that the greatest goods are the enjoyment of beautiful objects and personal affection, and that the existence of beauty unperceived has comparatively little value; our position is the same, and it suffices.

One may doubt whether beauty unseen is beauty at all. On AES-02, beauty is a response-dependent property, located in the relation between object and observer. A response-dependent property, however, is one that an object has in virtue of the response it would elicit in suitable observers, not the response it does elicit; a formation that no one will see is beautiful if it would elicit the response in a suitable observer, as AES-03's account of the convergent verdicts of ideal critics allows. The pair therefore does not require a mind-independent conception of beauty. It requires only that the dispositional property can be instantiated where no one exercises the disposition. Volume~V examines the case at length.

\subsection{Autonomy}

\paragraph{The pair.} Two outcomes are alike in the six other components. In the first, the members of a society choose their occupations, partners and convictions for themselves. In the second, a benevolent planner, knowing what each would choose and how each would fare, arranges matters so that each ends in the same occupation, with the same partner and the same convictions, through influences the members do not detect; the members' beliefs about how their lives came to be are false, and the loss of knowledge is compensated. The first outcome is better.

\paragraph{The reductions.} Mill's defense of liberty rests in part on the claim that people are the best judges of their own interests, so that self-direction tends to make lives go better \citep[ch.~III]{Mill1859}; on this reading the value of autonomy reduces to well-being. A second reduction treats autonomy as a form of achievement: a self-directed life is one whose shape is the agent's own work.

\paragraph{Why they fail.} The pair is constructed so that the planner's arrangement produces the same lives, the same experiences and the same distribution of the other goods as self-direction. If autonomy were valuable only as a means to well-being, the planner's arrangement, which secures the same well-being by other means, would be no worse. Mill also held that the free development of individuality is one of the elements of well-being, and \citet{Raz1986} that personal autonomy is valuable as a constituent of a good life, given a range of acceptable options; \citet{Dworkin1988} analyzes autonomy as a capacity to reflect on and endorse or revise one's first-order preferences. The pair shows what all three accounts need, that self-direction matters in itself. It also shows that autonomy is not achievement: the members' choices in the first outcome need not be difficult or competent, and the planner's arrangement does not deprive anyone of any outcome they aimed at, since each arrives where they would have aimed. What it deprives them of is the authorship of their lives. Whether autonomy is also a component of each person's well-being is examined in Volume~V; Chapter~7 of Volume~I argues that it is not wholly so, since the planner's subjects, judged by how their lives go for them, may fare exactly as well.

\begin{center}\small
\begin{tabular}{@{}p{2.2cm}p{4.6cm}p{4.4cm}r@{}}
\toprule
Component & Isolation pair (outline) & Principal reduction & Vol. \\
\midrule
Well-being & unremembered headache & (constituent view: no reduction) & II \\
Knowledge & accurate forgery; distant galaxy & to well-being; to achievement & II \\
Achievement & bridge fails by undetectable flaw & to knowledge and well-being & III \\
Relationships & friendship and its replica & to well-being; to knowledge & III \\
Fairness & consented lottery, same distribution & to equality, priority, desert & IV \\
Beauty & unseen formation & to aesthetic experience & V \\
Autonomy & benevolent planner & to well-being; to achievement & V \\
\bottomrule
\end{tabular}
\end{center}

\section{Composite and global reductions}\label{s:composite}

The reductions considered so far attempt to reduce one component to one or two others. Two more ambitious strategies remain.

\paragraph{Composite reductions.} A component might be reducible not to any single other component but to several jointly: achievement, for instance, to knowledge, autonomy and the well-being of those affected. Theorem~5.2 of ETH-01-a covers this case, since it concerns every set of components that omits the one in question, and a composite reduction of $i$ is a claim that $V$ supervenes on such a set. The isolation pair for achievement holds all six others fixed at once. Composite reductions therefore fail whenever the single-component reductions fail and the pairs are constructed, as ours are, with all six held fixed.

\paragraph{Global redescription.} A different strategy accepts that value does not supervene on any proper subset of the seven, and proposes a different set of dimensions altogether, on which value supervenes and which has fewer members. Corollary~\ref{c:unique} does not exclude this, since it concerns sets of the seven candidate components. The extreme case is a single dimension: value itself, or a magnitude defined from the seven levels by some function, such as \enquote{flourishing}. That value supervenes on such a dimension is trivial, since value supervenes on itself. But a monism of this kind is not a rival to the plurality thesis; it is a summary of it. The question that separates monists from pluralists is whether there is a single feature, specifiable without reference to the components, in virtue of which outcomes are better or worse. Pleasure, desire satisfaction and perfection have been the candidates, and each is either a part of a component, as pleasure is part of well-being, or a feature that idealized endorsement does not track when the components are held fixed. Corollary~5.3 of ETH-01-a, that value is not a function of any single component, together with the observation that every proposed single feature is either a component, a part of one, or definable from the components, exhausts the options known to us. The general claim, that no independent single feature could serve, is part of the completeness argument of Volume~VI, which identifies the closure of the list with the supervenience of value on all seven components together, and which therefore must show that any feature on which value supervenes is determined by them.

\section{Objections}\label{s:obj}

\begin{objection}[Entanglement]\label{o:entangle}
The components cannot be varied independently. There is no achievement without knowledge, no relationship without some well-being, no autonomy without some knowledge of one's options. Pairs that differ in only one component are therefore impossible, and the isolation test cannot be applied.
\end{objection}

\begin{reply}
The objection is right that the components are causally and sometimes constitutively connected, and wrong that the connection excludes the pairs. The test compares levels, and compensation allows the levels of the other components to be held fixed while the concrete outcomes differ (Section~\ref{s:method}). The engineer whose bridge fails has as much knowledge, at the level of the profile, as the engineer whose bridge stands, though not the same items of knowledge. Where compensation is in doubt, an alternative route is available: ETH-01-c derives the conclusion of ETH-01 from verdicts about pairs in which the outcome that is better in the tested component is no better, and may be worse, in the others, so that exact agreement in the other components is not required.
\end{reply}

\begin{objection}[Debunking]
The verdicts reflect the instrumental value of the components in ordinary circumstances. We are used to knowledge being useful and friendship making us happy, and our verdicts about cases in which they are not carry over the associations of ordinary cases.
\end{objection}

\begin{reply}
The objection describes a real source of error in judgments about isolated goods, but the idealized judge is not subject to it. She is fully informed about the outcomes, so she knows that the knowledge in the second pair for knowledge is useless, and she reasons without error, so she does not credit it with a usefulness she knows it lacks. The objection must therefore claim that our verdicts are not good estimates of hers. The comparative form of the pairs is designed to make them good estimates, in the way ETH-05-a shows for the experience machine: the debunking explanations of that case, from status quo bias to imaginative contamination, are each avoided by stating the case as a comparison of two outcomes rather than as a choice. The same form is used here. The burden is then on the debunker to identify, for each pair, an error that survives the comparative statement.
\end{reply}

\begin{objection}[Prudential goods only]
Relationships, knowledge, achievement and autonomy are goods for persons, and the impartial standpoint should therefore count them only through well-being. Listing them separately counts goods for persons twice: once as components of well-being, and once as components of the good.
\end{objection}

\begin{reply}
The isolation test does not allow double counting, since each component is tested with the others, including well-being, held fixed. If knowledge is also a constituent of well-being, its contribution to well-being is recorded in the well-being level, and the knowledge component records only the contribution that remains when that level is held fixed; the pairs show that such a contribution exists. Whether a good that passes the isolation test is a good for some person, a good of a relation between persons, or a good of no one in particular is a further question. Fairness is a good of no one in particular (Theorem~10.4 of Volume~IV); beauty, on the pair above, is too. For the other components the answer varies with the case, and the plurality thesis does not depend on it.
\end{reply}

\begin{objection}[Incomparability]
If the components are irreducible, there is no common measure of their value, and outcomes that trade one against another are incomparable. Pluralism of this kind leaves idealized endorsement silent in most real choices.
\end{objection}

\begin{reply}
Proposition~\ref{p:interact} shows that irreducibility does not imply incomparability: a value function can depend on all seven components, in no proper subset of them, and still rank every pair of profiles. Whether idealized endorsement in fact ranks every pair is a question about convergence, which ETH-01-a treats in Chapter~6 of Volume~I: where idealized agents diverge, there is no fact about which outcome is better, and many apparent cases of incomparability are cases of parity \citep{Chang2002}. That question is independent of the number of components, since a monistic theory can also face divergence in how its single good is measured.
\end{reply}

\section{What the discussion does not establish}\label{s:not}

First, the discussion does not establish that the list is complete. Proposition~\ref{p:essential} and Corollary~\ref{c:unique} show what follows if it is; the argument that it is belongs to Volume~VI, and the precise form of the thesis, supervenience of value on all seven components together, is stated in ETH-01-a and ETH-01-d.

Second, it does not establish the weights of the components or the form of $V$. The isolation verdicts show only that each component makes some difference. Proposition~\ref{p:interact} is an existence proof, not a proposal.

Third, it does not establish that idealized reflective endorsement converges. The criterion of Definition~3.2 of ETH-01-a is presupposed throughout. MTE-01 holds that moral facts are natural facts about idealized reactions, and the present discussion relies on that account rather than defending it.

Fourth, the pairs are stated in outline. Each depends on compensations whose adequacy is a matter for the volume devoted to the component, and a defect in one of them would affect the irreducibility of that component only, by Proposition~\ref{p:essential}, not the irreducibility of the others.

\section{Conclusion}

The irreducibility of a component of the good, in the sense ETH-01 asserts, is its membership in every supervenience base of value; given the completeness thesis, the seven components form the unique minimal base. Irreducibility constrains what value depends on, not how it depends on it, and so it is consistent with interaction among the components and with their full comparability. For each of the seven components there is a pair of outcomes, built where necessary by compensation within the other components, that idealized endorsement ranks differently, and each of the principal attempts to reduce a component to the others is refuted by its pair.

\begin{thebibliography}{99}\small
\bibitem[Bradford(2015)]{Bradford2015} Bradford, G. (2015). \emph{Achievement}. Oxford University Press.
\bibitem[Broome(1990)]{Broome1990} Broome, J. (1990). Fairness. \emph{Proceedings of the Aristotelian Society}, 91, 87--101.
\bibitem[Chang(2002)]{Chang2002} Chang, R. (2002). The possibility of parity. \emph{Ethics}, 112(4), 659--688.
\bibitem[De~Brigard(2010)]{DeBrigard2010} De~Brigard, F. (2010). If you like it, does it matter if it's real? \emph{Philosophical Psychology}, 23(1), 43--57.
\bibitem[Diamond(1967)]{Diamond1967} Diamond, P.~A. (1967). Cardinal welfare, individualistic ethics, and interpersonal comparison of utility: Comment. \emph{Journal of Political Economy}, 75(5), 765--766.
\bibitem[Dworkin(1988)]{Dworkin1988} Dworkin, G. (1988). \emph{The Theory and Practice of Autonomy}. Cambridge University Press.
\bibitem[Fletcher(2013)]{Fletcher2013} Fletcher, G. (2013). A fresh start for the objective-list theory of well-being. \emph{Utilitas}, 25(2), 206--220.
\bibitem[Greco(2010)]{Greco2010} Greco, J. (2010). \emph{Achieving Knowledge}. Cambridge University Press.
\bibitem[Griffin(1986)]{Griffin1986} Griffin, J. (1986). \emph{Well-Being: Its Meaning, Measurement and Moral Importance}. Clarendon Press.
\bibitem[Kagan(1998)]{Kagan1998} Kagan, S. (1998). Rethinking intrinsic value. \emph{The Journal of Ethics}, 2(4), 277--297.
\bibitem[Kolodny(2003)]{Kolodny2003} Kolodny, N. (2003). Love as valuing a relationship. \emph{The Philosophical Review}, 112(2), 135--189.
\bibitem[Korsgaard(1983)]{Korsgaard1983} Korsgaard, C.~M. (1983). Two distinctions in goodness. \emph{The Philosophical Review}, 92(2), 169--195.
\bibitem[Kvanvig(2003)]{Kvanvig2003} Kvanvig, J.~L. (2003). \emph{The Value of Knowledge and the Pursuit of Understanding}. Cambridge University Press.
\bibitem[Lackey(2007)]{Lackey2007} Lackey, J. (2007). Why we don't deserve credit for everything we know. \emph{Synthese}, 158(3), 345--361.
\bibitem[Mill(1859)]{Mill1859} Mill, J.~S. (1859). \emph{On Liberty}. John W. Parker and Son.
\bibitem[Moore(1903)]{Moore1903} Moore, G.~E. (1903). \emph{Principia Ethica}. Cambridge University Press.
\bibitem[Parfit(1984)]{Parfit1984} Parfit, D. (1984). \emph{Reasons and Persons}. Clarendon Press.
\bibitem[Parfit(1997)]{Parfit1997} Parfit, D. (1997). Equality and priority. \emph{Ratio}, 10(3), 202--221.
\bibitem[Rabinowicz and R{\o}nnow-Rasmussen(2000)]{RabinowiczRonnow2000} Rabinowicz, W., and R{\o}nnow-Rasmussen, T. (2000). A distinction in value: Intrinsic and for its own sake. \emph{Proceedings of the Aristotelian Society}, 100, 33--51.
\bibitem[Raz(1986)]{Raz1986} Raz, J. (1986). \emph{The Morality of Freedom}. Clarendon Press.
\bibitem[Rice(2013)]{Rice2013} Rice, C.~M. (2013). Defending the objective list theory of well-being. \emph{Ratio}, 26(2), 196--211.
\bibitem[Scheffler(1997)]{Scheffler1997} Scheffler, S. (1997). Relationships and responsibilities. \emph{Philosophy \& Public Affairs}, 26(3), 189--209.
\bibitem[Sidgwick(1907)]{Sidgwick1907} Sidgwick, H. (1907). \emph{The Methods of Ethics} (7th ed.). Macmillan.
\bibitem[Sosa(2007)]{Sosa2007} Sosa, E. (2007). \emph{A Virtue Epistemology: Apt Belief and Reflective Knowledge, Volume I}. Oxford University Press.
\bibitem[Sumner(1996)]{Sumner1996} Sumner, L.~W. (1996). \emph{Welfare, Happiness, and Ethics}. Clarendon Press.
\bibitem[Weijers(2014)]{Weijers2014} Weijers, D. (2014). Nozick's experience machine is dead, long live the experience machine! \emph{Philosophical Psychology}, 27(4), 513--535.
\end{thebibliography}

\end{document}
